\documentclass{ximera}
%nieuw 2 september 2026
\title{Vraagstukken elektrisch veld: conceptuele vragen}
\author{Daan Lipkens}

\begin{document}
\begin{abstract}
    Conceptuele vragen op het elektrisch veld.
\end{abstract}
\maketitle

%———————————————————————————————————————
% VRAAG 1: testlading verdubbelen
\begin{question}
    Je meet het elektrisch veld in een punt \textit{P} met een testlading $Q_t$. Vervolgens vervang je deze testlading door een nieuwe testlading met dubbel zo grote lading, $2Q_t$. Wat gebeurt er met de gemeten grootte van het elektrisch veld $E$ in punt \textit{P}?

    \begin{multipleChoice}
        \choice{$E$ wordt twee keer zo groot.}
        \choice{$E$ wordt twee keer zo klein.}
        \choice[correct]{$E$ blijft ongewijzigd.}
    \end{multipleChoice}

    \begin{hint}
        De kracht op de testlading verandert wel, maar wat gebeurt er als je die kracht deelt door de (grotere) testlading?
    \end{hint}

    \begin{solution}\nl
        Het elektrisch veld $E$ wordt volledig bepaald door de \textbf{bronlading}, niet door de testlading. Een grotere testlading ondervindt weliswaar een grotere kracht, maar deze kracht is exact evenredig groter, waardoor de verhouding $E = F_e/Q_t$ (en dus de gemeten veldsterkte) ongewijzigd blijft.
    \end{solution}
\end{question}
%———————————————————————————————————————

%———————————————————————————————————————
% VRAAG 2: nulpunt tussen 2 gelijke ladingen
\begin{question}
    Twee identieke positieve puntladingen bevinden zich op een onderlinge afstand $d$ van elkaar. Op welke plaats tussen de twee ladingen is het totale elektrisch veld gelijk aan nul? Verklaar waarom net op die plaats.

    \begin{hint}
        Beide ladingen zijn identiek en creëren dus, op gelijke afstand, een even groot veld.
    \end{hint}

    \begin{solution}\nl
        Exact in het \textbf{midden} tussen de twee ladingen, op een afstand $d/2$ van elk van beide. Omdat de ladingen identiek zijn en zich op gelijke afstand van dat punt bevinden, zijn de twee afzonderlijke velden even groot. Beide ladingen zijn bovendien positief, waardoor beide velden in dat middelpunt van elkaar wegwijzen (elk veld wijst immers weg van zijn eigen bronlading): de twee velden zijn dus exact tegengesteld gericht en heffen elkaar op.
    \end{solution}
\end{question}
%———————————————————————————————————————

%———————————————————————————————————————
% VRAAG 3: veldlijnen snijden elkaar niet
\begin{question}
    Leg uit waarom twee elektrische veldlijnen elkaar nooit kunnen snijden.

    \begin{hint}
        Wat stelt de richting van een veldlijn in een bepaald punt precies voor?
    \end{hint}

    \begin{solution}\nl
        In elk punt van de ruimte is er precies \textbf{één} elektrische veldvector $\vec{E}$: het totale veld in dat punt heeft een unieke grootte, richting en zin. Een veldlijn is raaklijnig aan die veldvector. Zouden twee veldlijnen elkaar in een punt snijden, dan zou dat betekenen dat het veld in dat punt tegelijk twee verschillende richtingen zou hebben, wat onmogelijk is.
    \end{solution}
\end{question}
%———————————————————————————————————————

%———————————————————————————————————————
% VRAAG 4: veld binnenin geleidende bol
\begin{question}
    Een massieve, neutrale metalen bol wordt in een uniform extern elektrisch veld geplaatst. Leg uit waarom het elektrisch veld binnenin de bol gelijk aan nul is, ook al is het externe veld dat niet.

    \begin{hint}
        Wat zou er gebeuren met de vrije elektronen in de bol als het veld binnenin niet nul zou zijn?
    \end{hint}

    \begin{solution}\nl
        De metalen bol bevat vrije elektronen. Zolang er binnenin de bol een elektrisch veld aanwezig is, ondervinden deze vrije elektronen een kracht en verplaatsen ze zich. Deze verplaatsing gaat door totdat het extra veld van de herverdeelde ladingen het externe veld precies compenseert: pas dan heerst er elektrostatisch evenwicht en stopt de verplaatsing van de elektronen. In dat evenwicht is het totale veld binnenin de geleider dus steeds nul.
    \end{solution}
\end{question}
%———————————————————————————————————————

%———————————————————————————————————————
% VRAAG 5: kooi van faraday vs. isolator
\begin{question}
    Zou een holle isolator (bijvoorbeeld een plastic doosje) even goed werken als afscherming tegen externe elektrische velden als een holle geleider? Verklaar.

    \begin{multipleChoice}
        \choice[correct]{Nee, want een isolator heeft geen vrije ladingen die zich kunnen herschikken om het externe veld te compenseren.}
        \choice{Ja, want elk materiaal polariseert in een elektrisch veld.}
        \choice{Ja, want isolatoren geleiden geen elektrische stroom.}
    \end{multipleChoice}

    \begin{hint}
        Elektrische schermwerking berust op de \textbf{verplaatsing van vrije ladingen}.
    \end{hint}

    \begin{solution}\nl
        Nee. Elektrische schermwerking (de kooi van Faraday) werkt doordat de \textbf{vrije elektronen} in een geleider zich herschikken tot het veld binnenin exact nul wordt. Een isolator bevat geen vrije ladingen: de atomen kunnen wel lichtjes polariseren, maar deze verschuiving is veel te klein om het externe veld volledig te compenseren. Een holle isolator biedt dus geen bruikbare afscherming.
    \end{solution}
\end{question}
%———————————————————————————————————————

%———————————————————————————————————————
% CONCEPT OEF VIERKANT 1: schaakbordpatroon, symmetrie geeft NUL veld
\begin{exercise}
    Vier puntladingen met dezelfde grootte $|Q|$ bevinden zich op de hoekpunten van een vierkant, zoals aangegeven op de tekening. Punt $P$ ligt in het midden van het vierkant.

    \begin{center}
        \begin{tikzpicture}[scale=0.9]
            \tkzDefPoint(0,4){A}
            \tkzDefPoint(4,4){B}
            \tkzDefPoint(4,0){C}
            \tkzDefPoint(0,0){D}
            \tkzDefPoint(2,2){P}
            \tkzDrawPolygon(A,B,C,D)
            \tkzLabelPoint[above=1.5mm](A){$+Q$}
            \tkzLabelPoint[above=1.5mm](B){$-Q$}
            \tkzLabelPoint[below=1.5mm](C){$+Q$}
            \tkzLabelPoint[below=1.5mm](D){$-Q$}
            \tkzLabelPoint[below=1.5mm](P){$P$}
            \draw[charge+] (A) circle (\R) node[scale=0.5]{$+$};
            \draw[charge-] (B) circle (\R) node[scale=0.5]{$-$};
            \draw[charge+] (C) circle (\R) node[scale=0.5]{$+$};
            \draw[charge-] (D) circle (\R) node[scale=0.5]{$-$};
            \filldraw (P) circle (1.5pt);
        \end{tikzpicture}
    \end{center}

    Naar waar wijst het resulterende elektrisch veld in punt $P$?

    \begin{multipleChoice}
        \choice{Naar boven.}
        \choice{Naar rechts.}
        \choice{Diagonaal, van de negatieve naar de positieve ladingen.}
        \choice[correct]{Het elektrisch veld is nul in $P$.}
    \end{multipleChoice}

    \begin{hint}
        Kijk naar de ladingen die \emph{diagonaal} tegenover elkaar liggen. Ze bevinden zich op dezelfde afstand van $P$ en hebben hetzelfde teken en dezelfde grootte.
    \end{hint}

    \begin{solution}\nl
        De ladingen op de twee uiteinden van eenzelfde diagonaal ($A$ en $C$, beide $+Q$; $B$ en $D$, beide $-Q$) liggen op exact dezelfde afstand van $P$ en hebben dezelfde grootte en hetzelfde teken. Het veld van $A$ in $P$ wijst dan precies tegengesteld aan het veld van $C$ in $P$, waardoor deze twee elkaar volledig opheffen. Hetzelfde geldt voor $B$ en $D$. Het resulterende veld in $P$ is dus \textbf{nul}.
    \end{solution}
\end{exercise}
%———————————————————————————————————————

%———————————————————————————————————————
% CONCEPT OEF VIERKANT 2: boven +, onder -, homogeen-achtig veld
\begin{exercise}
    Vier puntladingen met dezelfde grootte $|Q|$ bevinden zich op de hoekpunten van een rechthoek, zoals aangegeven op de tekening. Punt $P$ ligt in het midden van de rechthoek.

    \begin{center}
        \begin{tikzpicture}[scale=0.9]
            \tkzDefPoint(0,3){A}
            \tkzDefPoint(5,3){B}
            \tkzDefPoint(5,0){C}
            \tkzDefPoint(0,0){D}
            \tkzDefPoint(2.5,1.5){P}
            \tkzDrawPolygon(A,B,C,D)
            \tkzLabelPoint[above=1.5mm](A){$+Q$}
            \tkzLabelPoint[above=1.5mm](B){$+Q$}
            \tkzLabelPoint[below=1.5mm](C){$-Q$}
            \tkzLabelPoint[below=1.5mm](D){$-Q$}
            \tkzLabelPoint[right=1.5mm](P){$P$}
            \draw[charge+] (A) circle (\R) node[scale=0.5]{$+$};
            \draw[charge+] (B) circle (\R) node[scale=0.5]{$+$};
            \draw[charge-] (C) circle (\R) node[scale=0.5]{$-$};
            \draw[charge-] (D) circle (\R) node[scale=0.5]{$-$};
            \filldraw (P) circle (1.5pt);
        \end{tikzpicture}
    \end{center}

    Naar waar wijst het resulterende elektrisch veld in punt $P$?

    \begin{multipleChoice}
        \choice{Naar boven, richting de positieve ladingen.}
        \choice[correct]{Naar onder, richting de negatieve ladingen.}
        \choice{Naar links of naar rechts.}
        \choice{Het elektrisch veld is nul in $P$.}
    \end{multipleChoice}

    \begin{hint}
        Denk aan de horizontale en de verticale bijdrage van elke lading apart. Welke component (horizontaal of verticaal) valt weg door de symmetrie, en welke component versterkt net?
    \end{hint}

    \begin{solution}\nl
        Door de links-rechtssymmetrie van de figuur heffen de horizontale bijdragen van elk paar ladingen elkaar telkens op. Enkel de verticale bijdragen blijven over. De bovenste, positieve ladingen duwen $P$ weg van zichzelf, dus naar onder. De onderste, negatieve ladingen trekken $P$ naar zichzelf toe, eveneens naar onder. Alle verticale bijdragen werken dus in dezelfde zin en het resulterende veld in $P$ wijst \textbf{recht naar onder}, van de positieve naar de negatieve ladingen — vergelijkbaar met het homogene veld tussen twee geladen platen.
    \end{solution}
\end{exercise}
%———————————————————————————————————————

%———————————————————————————————————————
% CONCEPT OEF VIERKANT 3: ongelijke ladingen op de diagonaal
\begin{exercise}
    Vier puntladingen bevinden zich op de hoekpunten van een vierkant, zoals aangegeven op de tekening. Twee van de ladingen zijn gelijk in grootte en teken ($-Q$), maar op de andere diagonaal is één lading twee keer zo groot als de andere. Punt $P$ ligt in het midden van het vierkant.

    \begin{center}
        \begin{tikzpicture}[scale=0.9]
            \tkzDefPoint(0,4){A}
            \tkzDefPoint(4,4){B}
            \tkzDefPoint(4,0){C}
            \tkzDefPoint(0,0){D}
            \tkzDefPoint(2,2){P}
            \tkzDrawPolygon(A,B,C,D)
            \tkzLabelPoint[above=1.5mm](A){$+2Q$}
            \tkzLabelPoint[above=1.5mm](B){$-Q$}
            \tkzLabelPoint[below=1.5mm](C){$+Q$}
            \tkzLabelPoint[below=1.5mm](D){$-Q$}
            \tkzLabelPoint[below=1.5mm](P){$P$}
            \draw[charge+] (A) circle (\R) node[scale=0.5]{$+$};
            \draw[charge-] (B) circle (\R) node[scale=0.5]{$-$};
            \draw[charge+] (C) circle (\R) node[scale=0.5]{$+$};
            \draw[charge-] (D) circle (\R) node[scale=0.5]{$-$};
            \filldraw (P) circle (1.5pt);
        \end{tikzpicture}
    \end{center}

    Naar waar wijst het resulterende elektrisch veld in punt $P$?

    \begin{multipleChoice}
        \choice{Diagonaal, richting $A$ (de sterkste lading).}
        \choice[correct]{Diagonaal, richting $C$ (weg van de sterkste lading).}
        \choice{Horizontaal, van $D$ naar $C$.}
        \choice{Het elektrisch veld is nul in $P$.}
    \end{multipleChoice}

    \begin{hint}
        Ladingen $B$ en $D$ heffen elkaar volledig op (zelfde grootte, zelfde teken, diagonaal tegenover elkaar). Wat blijft er dan nog over?
    \end{hint}

    \begin{solution}\nl
        De ladingen $B$ en $D$ zijn identiek ($-Q$) en liggen diagonaal tegenover elkaar op dezelfde afstand van $P$: hun bijdragen heffen elkaar volledig op, net als in de eerste oefening van deze reeks.

        Wat overblijft, is de bijdrage van $A$ ($+2Q$) en $C$ ($+Q$), eveneens diagonaal tegenover elkaar op gelijke afstand van $P$, maar nu met een \emph{ongelijke} grootte. Beide zijn positief, dus duwen ze $P$ weg van zichzelf. Omdat $A$ twee keer zo sterk is als $C$, wint de bijdrage van $A$ het pleit. Het resulterende veld in $P$ wijst dus diagonaal, \textbf{weg van $A$ en richting $C$}.
    \end{solution}
\end{exercise}
%———————————————————————————————————————

%———————————————————————————————————————
% CONCEPT OEF 2 LADINGEN 1: tegengesteld teken, GELIJKE grootte (dipool) -> geen nulpunt
\begin{exercise}
    Twee tegengestelde puntladingen $+Q$ en $-Q$, met exact dezelfde grootte, bevinden zich op een rechte lijn.

    \begin{center}
        \begin{tikzpicture}[scale=1]
            \tkzDefPoint(0,0){Q1}
            \tkzDefPoint(6,0){Q2}
            \draw[gray,thick] (Q1) -- (Q2);
            \tkzLabelPoint[below=1.5mm](Q1){$+Q$}
            \tkzLabelPoint[below=1.5mm](Q2){$-Q$}
            \draw[charge+] (Q1) circle (\R) node[scale=0.5]{$+$};
            \draw[charge-] (Q2) circle (\R) node[scale=0.5]{$-$};
        \end{tikzpicture}
    \end{center}

    Is er op de lijn door beide ladingen een plaats waar het elektrisch veld nul is? Verklaar.

    \begin{multipleChoice}
        \choice{Ja, ergens tussen de twee ladingen.}
        \choice{Ja, links van $+Q$.}
        \choice{Ja, rechts van $-Q$.}
        \choice[correct]{Nee, nergens op een eindige afstand (enkel op oneindige afstand nadert het veld nul).}
    \end{multipleChoice}

    \begin{hint}
        Ga voor elk van de drie gebieden (links, tussen, rechts) na of de twee velden elkaar ergens tegenwerken.
    \end{hint}
    \begin{hint}
        Vergelijk deze situatie met de vorige oefening: wat was daar de reden dat er wél een nulpunt bestond? Is die reden hier ook aanwezig?
    \end{hint}

    \begin{solution}\nl
        \emph{Tussen} de twee ladingen wijzen beide velden in dezelfde zin (weg van $+Q$, richting $-Q$): ze werken elkaar dus nooit tegen en tellen enkel op.

        \emph{Links van $+Q$ of rechts van $-Q$} wijzen de twee velden wel in tegengestelde zin. Maar omdat beide ladingen exact even groot zijn, is de lading die het dichtst bij ligt altijd dominant: op elke eindige afstand is de bijdrage van de dichtstbijzijnde lading sterker dan die van de verder gelegen lading (met dezelfde grootte). De twee velden kunnen elkaar dus nooit precies opheffen.

        In tegenstelling tot de vorige oefening (ongelijke ladingen) bestaat hier dus \textbf{geen enkel punt op eindige afstand} waar het elektrisch veld nul is. Enkel oneindig ver van beide ladingen nadert het veld naar nul.
    \end{solution}
\end{exercise}
%———————————————————————————————————————

%———————————————————————————————————————
% CONCEPT OEF 2 LADINGEN 2: zelfde teken, ongelijke grootte -> nulpunt TUSSEN
\begin{exercise}
    Twee positieve puntladingen $Q_{1}$ en $Q_{2}$ bevinden zich op een rechte lijn, met $Q_{1} < Q_{2}$ (dus $Q_2$ is de sterkste lading).

    \begin{center}
        \begin{tikzpicture}[scale=1]
            \tkzDefPoint(0,0){Q1}
            \tkzDefPoint(6,0){Q2}
            \draw[gray,thick] (Q1) -- (Q2);
            \tkzLabelPoint[below=1.5mm](Q1){$Q_1$}
            \tkzLabelPoint[below=1.5mm](Q2){$Q_2$}
            \draw[charge+] (Q1) circle (\R) node[scale=0.5]{$+$};
            \draw[charge+] (Q2) circle (\R) node[scale=0.6]{$+$};
        \end{tikzpicture}
    \end{center}

    Is er op de lijn door beide ladingen een plaats waar het elektrisch veld nul is? Zo ja, waar ligt die plaats ongeveer?

    \begin{multipleChoice}
        \choice{Nee, er is nergens een punt waar het veld nul is.}
        \choice{Ja, tussen de ladingen, dichter bij de sterkste lading $Q_2$.}
        \choice[correct]{Ja, tussen de ladingen, dichter bij de zwakste lading $Q_1$.}
        \choice{Ja, maar buiten het gebied tussen de twee ladingen, het korste bij lading $Q_1$.}
        \choice{Ja, maar buiten het gebied tussen de twee ladingen, het korste bij lading $Q_2$.}
    \end{multipleChoice}

    \begin{hint}
        Bekijk eerst de zin van het veld in de drie gebieden (links van $Q_1$, tussen $Q_1$ en $Q_2$, rechts van $Q_2$). In welk gebied wijzen de twee bijdragen elkaar tegen?
    \end{hint}
    \begin{hint}
        Om dezelfde veldsterkte te bekomen met een kleinere lading, moet je dichter bij die lading staan.
    \end{hint}

    \begin{solution}\nl
        Omdat beide ladingen positief zijn, wijzen hun velden in het gebied \emph{tussen} $Q_1$ en $Q_2$ in tegengestelde zin naar elkaar toe (weg van elk van beide ladingen). Links van $Q_1$ of rechts van $Q_2$ wijzen beide velden juist in dezelfde zin en kunnen ze nooit nul worden. Het nulpunt ligt dus ergens tussen de twee ladingen.

        Om het veld van de zwakkere lading $Q_1$ even groot te maken als dat van de sterkere lading $Q_2$, moet je dichter bij $Q_1$ staan (kleinere afstand compenseert de kleinere lading). Het nulpunt ligt dus tussen de ladingen, \textbf{dichter bij de zwakste lading $Q_1$}.
    \end{solution}
\end{exercise}
%———————————————————————————————————————

%———————————————————————————————————————
% CONCEPT OEF 2 LADINGEN 3: verschillend teken, ongelijke grootte -> nulpunt TUSSEN
\begin{exercise}
    Twee puntladingen $Q_{1}$ en $Q_{2}$ bevinden zich op een rechte lijn, met $Q_{1} < 0$ en $Q_{2} > 0$, en $|Q_{1}| < |Q_{2}|$ (dus $Q_1$ is de zwakkere negatieve lading en $Q_2$ is de sterkere positieve lading).

    \begin{center}
        \begin{tikzpicture}[scale=1]
            \tkzDefPoint(0,0){Q1}
            \tkzDefPoint(6,0){Q2}
            \draw[gray,thick] (Q1) -- (Q2);
            \tkzLabelPoint[below=1.5mm](Q1){$Q_1$}
            \tkzLabelPoint[below=1.5mm](Q2){$Q_2$}
            \draw[charge-] (Q1) circle (\R) node[scale=0.5]{$-$};
            \draw[charge+] (Q2) circle (\R) node[scale=0.6]{$+$};
        \end{tikzpicture}
    \end{center}

    Is er op de lijn door beide ladingen een plaats waar het elektrisch veld nul is? Zo ja, waar ligt die plaats ongeveer?

    \begin{multipleChoice}
        \choice{Nee, er is nergens een punt waar het veld nul is.}
        \choice{Ja, tussen de ladingen, dichter bij de sterkste lading $Q_2$.}
        \choice{Ja, tussen de ladingen, dichter bij de zwakste lading $Q_1$.}
        \choice[correct]{Ja, maar buiten het gebied tussen de twee ladingen, het korste bij lading $Q_1$.}
        \choice{Ja, maar buiten het gebied tussen de twee ladingen, het korste bij lading $Q_2$.}
    \end{multipleChoice}

    \begin{hint}
        Bekijk eerst de zin van het veld in de drie gebieden (links van $Q_1$, tussen $Q_1$ en $Q_2$, rechts van $Q_2$). In welk gebied wijzen de twee bijdragen elkaar tegen?
    \end{hint}
    \begin{hint}
        Om dezelfde veldsterkte te bekomen met een kleinere lading, moet je dichter bij die lading staan.
    \end{hint}

    \begin{solution}\nl
        Omdat beide ladingen verschillend zijn, wijzen hun velden in het gebied \emph{buiten} $Q_1$ en $Q_2$ in tegengestelde zin naar elkaar toe (weg van de positieve lading $Q_2$ en naar de negatieve lading $Q_1$).

        Om het veld van de zwakkere lading $Q_1$ even groot te maken als dat van de sterkere lading $Q_2$, moet je dichter bij $Q_1$ staan (kleinere afstand compenseert de kleinere lading). Het nulpunt ligt dus links van lading $Q_1$, of dus buiten de twee ladingen het korste bij de kleinse lading $Q_1$.
    \end{solution}
\end{exercise}
%———————————————————————————————————————

\end{document}